\scanpage{44}
% L1
par
\[
\mathrm{cl}(X)\mapsto\sum(-1)^i\,\mathrm{cl}\,H^i_!(X,\mathbb Q(0))
\]
où bien entendu $H^i_!(X,\mathbb Z(0))$ est \deleted{\unsure{considéré}}\inserted{défini} comme motif sur $K$ par $R^if(\mathbb Q(0))$ ; $f:X\mathrel{\unsure{-|\to}}\operatorname{Spec}K$ la projection.

% L6
Prenant les composantes homogènes de $\sigma\mathrm{cl}(X)$, on a des $\sigma_i(\mathrm{cl}(X))$ qui sont des ``motifs virtuels'', et permettent de définir des ``nbs de Betti virtuels'', ceux demandés par Serre.

% L11 — diagonal crosses through the following paragraphs
\deleted{Si on filtre $L(K)$ par la dimension en des $L_i(K)$, on trouve}
\[
\deleted{L_i(K)\longrightarrow D_i(M^+_{\unsure{\xi|\varepsilon}}(K))}
\]
\deleted{On peut se demander si c'est surjectif, ou injectif, \deleted{ou bijectif mod} éventuellement mod groupes de torsion. Déjà pour $i=1$ ($K$ alg clos \inserted{de degré de trans. infini sur le corps premier}) rien n'est connu.}
% slanted
\margin{Pb pas bien posé \unsure{car avons|car on a} réponse \uline{négative} en \deleted{\unsure{telle}} \unsure{partie|partout}. cf plus bas la formulation correcte \dots\ En fait $L_i(K)\to M^+_{(i)}(K)+\xi M^+_{(i-1)}(K)+\xi^2 M^+_{(i-2)}(K)+\dotsb$}
\deleted{On peut corser la situation en considérant des $X/K$ munis d'un grpe fini donné $G$ d'autom. -- mais il faut}
% archive sheet number: 43
